% Hafta 5 — XXE: XML dış varlık saldırısı
% Derleme: py -3.12 tools/latex/derle.py docs/week-5/assets/h05-07-xxe.tex
\documentclass[tikz,border=8pt]{standalone}
\usepackage{cen429-cizim}
\begin{document}
\begin{tikzpicture}
  \node[baslik] (b) at (0,0) {XXE: XML dış varlık saldırısı};
  \node[altbaslik] at (b.south west) {ayrıştırıcının ``yardımsever'' varsayılanı};
  \node[kotukutu=38mm, anchor=north west] (n1) at (0,-2.4)
    {\kb{XML girdisi}\\\kod{<!DOCTYPE ... SYSTEM}\\\kod{"file:///etc/passwd">}};
  \node[uyarikutu=32mm, right=9mm of n1] (n2)
    {\kb{Ayrıştırıcı}\\dış varlığı ÇÖZER};
  \node[kotukutu=32mm, right=9mm of n2] (n3)
    {\kb{Dosya okunur}\\sunucudan veri sızar};
  \node[iyikutu=38mm, right=9mm of n3] (n4)
    {\kb{Kapatan ayar}\\disallow-doctype-decl\\external entities = false};
  \draw[kotuok] (n1.east) -- (n2.west);
  \draw[kotuok] (n2.east) -- (n3.west);
  \draw[iyiok] (n3.east) -- (n4.west);
  \node[fit=(n1)(n2)(n3)(n4), inner sep=0pt] (grp) {};
  \node[dip, text width=155mm, anchor=north west] at ($(grp.south west)+(0,-8mm)$)
    {Kural: XML/YAML/şablon ayrıştırıcılarında gereksiz yetenekleri KAPATIN.};
\end{tikzpicture}
\end{document}
